\documentclass[11pt]{article}

% ---------- Page ----------
\usepackage[margin=1in]{geometry}

% ---------- Fonts (XeLaTeX/LuaLaTeX only) ----------
\usepackage{fontspec}

\usepackage[T1]{fontenc}
\usepackage{lmodern}   % clean, modern LaTeX font


% ---------- Typography ----------
\usepackage{microtype}
\usepackage{setspace}
\setstretch{1.08}
\usepackage{inconsolata}
% ---------- Links ----------
\usepackage[hidelinks]{hyperref}

% ---------- Colors ----------
\usepackage{xcolor}
\definecolor{ink}{HTML}{111827}      % slate-900
\definecolor{sub}{HTML}{374151}      % slate-700
\definecolor{muted}{HTML}{6B7280}    % slate-500
\definecolor{line}{HTML}{E5E7EB}     % slate-200
\definecolor{panel}{HTML}{F9FAFB}    % slate-50
\definecolor{accent}{HTML}{2563EB}   % blue-600
\definecolor{codebg}{HTML}{111827}   % dark slate for code background

% ---------- Headings ----------
\usepackage{titlesec}
\titleformat{\section}
  {\sffamily\Large\bfseries\color{ink}}{\thesection}{0.75em}{}
\titleformat{\subsection}
  {\sffamily\large\bfseries\color{ink}}{\thesubsection}{0.75em}{}
\titleformat{\subsubsection}
  {\sffamily\normalsize\bfseries\color{ink}}{\thesubsubsection}{0.75em}{}

% ---------- Lists ----------
\usepackage{enumitem}
\setlist[itemize]{topsep=4pt,itemsep=2pt,leftmargin=*}
\setlist[enumerate]{topsep=4pt,itemsep=2pt,leftmargin=*}

% ---------- Modern boxes ----------
\usepackage[most]{tcolorbox}
\tcbset{
  enhanced,
  boxrule=0.6pt,
  colframe=line,
  colback=panel,
  arc=8pt,
  left=10pt,right=10pt,top=8pt,bottom=8pt
}

% ---------- Code (minted only) ----------
\usepackage{minted}
\setminted{
  style=monokai,
  bgcolor=codebg,
  fontsize=\small,
  linenos,
  breaklines,
  autogobble
}

% ---------- Custom tcolorbox wrappers ----------
\newtcolorbox{cmd}[1][]{%
  title={\sffamily\bfseries Command%
    \if\relax\detokenize{#1}\relax\else\ — #1\fi},
  coltitle=ink,
  colback=panel,
  colframe=line
}

\newtcolorbox{out}[1][]{%
  title={\sffamily\bfseries Output%
    \if\relax\detokenize{#1}\relax\else\ — #1\fi},
  coltitle=ink,
  colback=panel,
  colframe=line
}

% ---------- Header / Footer ----------
\usepackage{fancyhdr}
\pagestyle{fancy}
\fancyhf{}
\setlength{\headheight}{14pt}
\lhead{\textcolor{muted}{MongoDB Indexing Lab}}
\rhead{\textcolor{muted}{\today}}
\cfoot{\textcolor{muted}{\thepage}}

% ---------- Title ----------
\usepackage{titling}
\pretitle{\vspace{-1.0cm}\color{ink}\LARGE\sffamily\bfseries}
\posttitle{\par\vspace{0.15cm}\color{sub}\large\sffamily}
\preauthor{\color{sub}\normalsize\sffamily}
\postauthor{\par}
\predate{\color{muted}\normalsize\sffamily}
\postdate{\par\vspace{0.6cm}\hrule\vspace{0.8cm}}

\title{Indexing and Query Performance in MongoDB}
\author{Nuria Arroyo \quad \texttt{(173605)}}
\date{MongoDB Compass + Atlas \; \textbullet \; Database \texttt{tp1}, Collection \texttt{restaurantes}}


\begin{document}
\maketitle

\begin{tcolorbox}
\textbf{Objective.} Explore the impact of indexing on query performance in MongoDB. We run baseline queries without indexes, inspect execution plans using \texttt{explain}, create the required indexes, re-run the same queries, and compare performance metrics. Finally, we drop the indexes and observe the regression.
\end{tcolorbox}

\section{Environment}
\begin{itemize}
  \item Platform: MongoDB Atlas (cloud) + MongoDB Compass (mongosh inside Compass)
  \item Database: \texttt{tp1}
  \item Collection: \texttt{restaurantes}
  \item Dataset size (approx.): \underline{\hspace{3cm}} documents
\end{itemize}

\section{Procedure (Assignment Instructions)}
\begin{enumerate}
  \item Execute queries \textbf{without indexes} and observe performance using \texttt{explain}.
  \item Create the specified indexes (Section 3.5).
  \item View / verify indexes.
  \item Execute the same queries \textbf{with indexes} and observe changes.
  \item Compare and analyze results.
  \item Delete indexes and re-check query performance.
\end{enumerate}
\newpage

\section{Step 1 — Baseline: queries without indexes}
\subsection{Confirm existing indexes}
\begin{cmd}[Verify index state (baseline)]
\begin{minted}{javascript}
use tp1
db.restaurantes.getIndexes()
\end{minted}
\end{cmd}

\begin{out}[Paste your output here]
\begin{minted}[linenos=false]{text}
% Example:
% [ { v: 2, key: { _id: 1 }, name: '_id_' } ]
\end{minted}
\end{out}

\subsection{Baseline query: all Italian restaurants}
\begin{cmd}[Explain with execution stats (baseline)]
\begin{minted}{javascript}
db.restaurantes.find({ cuisine: "Italian" }).explain("executionStats")
\end{minted}
\end{cmd}

\begin{out}[Paste your explain output here]
\begin{minted}[linenos=false]{text}
% Paste the full explain("executionStats") output here
\end{minted}
\end{out}

\subsection{Baseline notes}
From the explain output, record:
\begin{itemize}
  \item Winning plan stage (expected: \texttt{COLLSCAN} without an index on \texttt{cuisine})
  \item \texttt{nReturned}
  \item \texttt{totalDocsExamined}
  \item \texttt{totalKeysExamined}
  \item \texttt{executionTimeMillis}
\end{itemize}

\begin{tcolorbox}
\textbf{Baseline summary (fill in).}\\
Winning plan: \underline{\hspace{5.5cm}} \\
nReturned: \underline{\hspace{3.5cm}} \quad
Docs examined: \underline{\hspace{3.5cm}} \\
Keys examined: \underline{\hspace{3.5cm}} \quad
Time (ms): \underline{\hspace{3.5cm}}
\end{tcolorbox}

\section{Step 2 — Create indexes (Section 3.5)}
\subsection{Q20: increasing index on \texttt{cuisine}}
\begin{cmd}[Create single-field index]
\begin{minted}{javascript}
db.restaurantes.createIndex({ cuisine: 1 })
\end{minted}
\end{cmd}

\begin{out}[Paste output (index name)]
\begin{minted}[linenos=false]{text}
% Example:
% cuisine_1
\end{minted}
\end{out}

\subsection{Q21: compound index (\texttt{cuisine} asc, \texttt{address.zipcode} desc)}
\begin{cmd}[Create compound index]
\begin{minted}{javascript}
db.restaurantes.createIndex({ cuisine: 1, "address.zipcode": -1 })
\end{minted}
\end{cmd}

\begin{out}[Paste output (index name)]
\begin{minted}[linenos=false]{text}
% Example:
% cuisine_1_address.zipcode_-1
\end{minted}
\end{out}

\section{Step 3 — View indexes}
\begin{cmd}[Verify indexes created]
\begin{minted}{javascript}
db.restaurantes.getIndexes()
\end{minted}
\end{cmd}

\begin{out}[Paste output]
\begin{minted}[linenos=false]{text}
% Paste the output here
\end{minted}
\end{out}

\section{Step 4 — Re-run queries with indexes}
\subsection{Italian restaurants query (with indexes)}
\begin{cmd}[Explain with execution stats (indexed)]
\begin{minted}{javascript}
db.restaurantes.find({ cuisine: "Italian" }).explain("executionStats")
\end{minted}
\end{cmd}

\begin{out}[Paste explain output]
\begin{minted}[linenos=false]{text}
% Paste the full explain("executionStats") output here
\end{minted}
\end{out}

\begin{tcolorbox}
\textbf{Indexed summary (fill in).}\\
Winning plan: \underline{\hspace{5.5cm}} \\
nReturned: \underline{\hspace{3.5cm}} \quad
Docs examined: \underline{\hspace{3.5cm}} \\
Keys examined: \underline{\hspace{3.5cm}} \quad
Time (ms): \underline{\hspace{3.5cm}}
\end{tcolorbox}

\section{Step 5 — Analysis}
\subsection{Comparison table}
\begin{tcolorbox}
\begin{tabular}{p{4.2cm}p{5.0cm}p{5.0cm}}
\textbf{Metric} & \textbf{Baseline (no index)} & \textbf{With index} \\
\hline
Winning plan stage & \underline{\hspace{4.5cm}} & \underline{\hspace{4.5cm}} \\
nReturned & \underline{\hspace{4.5cm}} & \underline{\hspace{4.5cm}} \\
Docs examined & \underline{\hspace{4.5cm}} & \underline{\hspace{4.5cm}} \\
Keys examined & \underline{\hspace{4.5cm}} & \underline{\hspace{4.5cm}} \\
Execution time (ms) & \underline{\hspace{4.5cm}} & \underline{\hspace{4.5cm}} \\
\end{tabular}
\end{tcolorbox}

\subsection{Discussion}
Write 4--8 sentences interpreting the differences. Mention whether the plan changed from
\texttt{COLLSCAN} to \texttt{IXSCAN}, and how \texttt{totalDocsExamined} and time changed.

\vspace{2.2cm}

\section{Step 6 — Delete indexes and re-check}
\subsection{Drop created indexes}
\begin{cmd}[Drop indexes (replace names with your actual index names)]
\begin{minted}{javascript}
db.restaurantes.dropIndex("cuisine_1")
db.restaurantes.dropIndex("cuisine_1_address.zipcode_-1")
\end{minted}
\end{cmd}

\begin{out}[Paste output]
\begin{minted}[linenos=false]{text}
% Paste output of dropIndex commands here
\end{minted}
\end{out}

\subsection{Explain again after dropping}
\begin{cmd}[Explain after dropping indexes]
\begin{minted}{javascript}
db.restaurantes.find({ cuisine: "Italian" }).explain("executionStats")
\end{minted}
\end{cmd}

\begin{out}[Paste explain output]
\begin{minted}[linenos=false]{text}
% Paste output here
\end{minted}
\end{out}

\subsection{Post-drop discussion}
Write 2--5 sentences describing what changed (plan stage, docs examined, time).

\vspace{2.0cm}

\end{document}